\section{From Crisis to Revolt}

\begin{quotex}
I don't know if the following review by \textbf{Rene Guenon} of Revolt against the Modern World has been published anywhere, but it is worth including here for the issues it raises. Like \textbf{Ananda Coomaraswamy}, Guenon's primary objection is also the emphasis of regality over the sacerdotal caste. Those who would take Evola's position seem to do so because they object to the word ``supremacy'', and then interpret the relationship between priest and king as a one way, rather than a mutual, relationship.
\end{quotex}


Evola's case is weakened by his misunderstanding of the rite of the purohita. Evola simply pulled the relevant reference to the Aitareya Brahmana from subsequent editions of the book, but continued the same theme as though it had never existed. This makes Evola's discussion on that topic suspect; besides, he offers many historical examples, but is short on demonstrating the principles that would support his case.

In his review of Coomaraswamy's \textit{Spiritual Authority and Temporal Power in the Indian Theory of Government}, Guenon delves deeper into this topic, from the point of view of principle. As Mr. Shankar has pointed out in a recent comment, this needs to start with the principle that unites priest and king rather than on what divides them. Furthermore, as the exterior is a reflection of the interior, we need to discern the roles of priest and king in our inner life.

First of all, Evola misrepresents the role of King in the \textbf{Ancient City} as documented by \textbf{Fustel de Coulanges}. It wasn't the case that the King had priority; rather that the roles of both Priest and King were held by one man, and thus were united. The regression of castes, then, occurred with the rebellion of the aristocracy, who split up the roles of the Priest-King, stripping him of his kingly function, leaving him as the High Priest of the city.

In that case, the understanding of the relationship between the spiritual and temporal makes little sense in Evola's scheme. The King's role is always in the temporal realm, or manifestation, which is the reflection of the spiritual, and the unmanifest, the province of the Priest. Again, this is necessarily a mutual relationship, so the question of supremacy is irrelevant. As for the King standing in for God, as Evola claims, this is a partial picture. In the Medieval system, Christ is both Priest and King. Hence, the King is God ordained, but that does not exhaust that claim. For the High Priest is likewise the Vicar of Christ in his own domain.

The King rules by edict, that is, the word; spiritual authority, however, comes from the idea, or thought\footnote{\url{https://gornahoor.net/?p=3936}}, to which the word is the reflection on the plane of manifestation. Hence, the King lays down the law, which is supreme in his realm, but infallibility has never been claimed by the King; that is a quality of the Priest.

Guenon points out that, viewed from the microcosm, these roles appear as the two atmas, i.e., the Self and the Ego. Establishing their proper relationship within, is part of the spiritual task, or ``Holy War''. Coomarawamy writes:

\begin{quotex}
The outer, active, feminine, and mortal aspect of our nature subsists more eminently in its inner, contemplative, masculine, and immortal side, to which it can and must be `reduced', that is to say, brought back or reunited.
\end{quotex}


Guenon comments:

\begin{quotex}
For a king, autonomy consists in not letting himself be ruled by the multitude subordinate to him, and likewise for each person, in not letting himself by ruled by the inferior and contingent elements of his being. Hence the two meanings of the `holy war' for the establishment and maintenance of order in both cases.
\end{quotex}


By fighting this war within, the Priest and King are united in a common battle. As Coomarawamy pointed out in his review of Revolt Against the Modern World\footnote{\url{https://gornahoor.net/?p=4274}}, Evola tends to absolutize this issue, categorizing things in absolute dichotomies of masculine/feminine or solar/lunar. However, these are relative terms. The King is masculine in respect to his realm and his subjects, but feminine to the spiritual and transcendental. To deny this is to make the King something he clearly is not; it also has deleterious consequences to one's inner spiritual life.

The following review was originally published in \emph{Le Voile d'Isis}, Paris, May 1934.

\begin{quotex}
In this new work, the author opposes traditional civilization to modern civilization, the former having a transcendent and essentially hierarchical character, and the latter founded on a purely human and contingent element. He then describes the phases of spiritual decadence that has led from the traditional world to the modern world. We will have some reservations to make on some points: so, when dealing with the original source of the two sacerdotal and regal powers, the author has a very definite tendency to put the emphasis on the regal aspect at the expense of the sacerdotal aspect. When he distinguishes two types of tradition that he relates respectively to the North and the South, the latter of these two terms appears to us as a little inaccurate, even if he does not mean it in a strictly ``geographic'' sense, since it seems to pertain especially to Atlantis, which, in every way, corresponds to the West, not the South. We also fear that he sees in primitive Buddhism something different from what it was in reality, since he praises it in way that is not absolutely encompassed in the traditional point of view. On the other hand, he disdains Pythagoreanism in a manner that is poorly justified. We could also point out other similar things.

This must not prevent us from recognizing, as is right, the merit and interest of the work as a whole, and to bring it in a particular way to the attention of all those who are concerned with the ``crisis of the modern world'', and who think like us that the only effective means of remediating it would consist in a return to the traditional spirit beyond which nothing truly ``constructive'' could be effectively undertaken.
\end{quotex}



\flrightit{Posted on 2012-05-29 by Cologero}

\begin{center}* * *\end{center}

\begin{footnotesize}
\begin{sffamily}

\texttt{Angolmois on 2012-05-31 at 01:50 said:}

Evola's mistakes and mis-representations are too serious not to be pin-pointed and repeated. As in the case of Nietzsche as well, his personal life is the best proof that there were some black holes in his thought and ``philosophy''. I see no ``obsession'' in Cologero's points, but I do see a fundamental mis-representation and mis-understanding of Evola in many other mansions of the anglosphere.

\hfill

\texttt{Mihai on 2012-05-31 at 05:47 said:}

@HOO: I am afraid that when discussing Evola, things tend to degenerate from rational debate to name-calling and one-sided blind arguments rather quickly, sometimes resembling arguments fought over football teams.

The problem is that it is nearly impossible to completely exclude the individual equation as long as one is still human, so a degree of partiality, no matter how small, is normal- in all these thinkers.

My take on this subject is like this: a king, as long as he embodies both spiritual authority and temporal power- that is he exercises both vertical and horizontal power- is an image of the Divine, of the Center. However, the spiritual aspect in him- the vertical dimension- is superior to his warrior aspect- the horizontal one. As a function, the king is the synthesis of both priest and warrior and as such he rules over both.

Where Evola makes a mistake is where he tries to make these two aspects- warrior and sacerdotal- as being equal, with no importance in which of these holds the higher place.

Furthermore, I think this site is quite respectful and well inclined towards Evola, the only place where it has reservations are some of his more controversial aspects.

\hfill

\texttt{Angolmois on 2012-05-31 at 11:02 said:}

@HOO: Guénon's own mistakes or mis-understadings (of Buddhism, towards which he changed his attitude after Coomaraswamy) do not relate to his right understanding of Evola's mis-representation of the relationship between royalty and pontificate. I personally agree with Mihai's understanding of the issue.

In discussing `The Hermetic Tradition', I think it is good to understand that the royal path of alchemy is more cosmic and psychic than universal and ``beyond cosmos''. This makes it suitable for hermeticists and for the royal caste, and the relation that Evola makes between `the white phase' and `the red phase', making the latter superior to the former, is right only in its own context. Yet, tt should not be made a general principle, since it relates to the supra-individual concept of `coming back to earth' (descending realization) and not into the relationship between sacerdotal caste and the royal caste.

\hfill

\end{sffamily}
\end{footnotesize}

